\chapter{正交集与线性方程组}

\section{本章导引}

本章专门讨论线性方程组.在处理线性方程组时,最常见的问题包括：
\begin{itemize}
	\item 确定线性方程组是否相容(是否至少存在一个解).
	\item 确定线性方程组是否具有唯一解.
	\item 寻找通解(如果存在).
	\item 检验两个方程组是否等价(是否具有相同的解).
	\item 求解关于某些选定变量的方程组.
	\item 已知初始线性方程组解的情况下,求解由于以下改变而产生的新方程组的解：
	\begin{itemize}
		\item 修改了一个或多个方程;
		\item 删除了一个或多个方程和(或)变量;
		\item 增加了一个或多个方程和(或)变量.
	\end{itemize}
\end{itemize}

在第2节中,我们分析相容性问题,并展示了如何以一种优雅的方式利用正交集来解决这些问题.

在第3节中,我们利用第1章中的正交化技术求解线性方程组的问题,并讨论了解的唯一性问题.

在第4节中,我们计算了所提方法的复杂度,并将其与Gauss消元法进行了对比.

在第5节中,我们应用该方法来分析两个方程组的等价性.在第6节中,我们展示了如何应用所提方法来求解关于某些选定变量的方程组.

大多数现有的求解方程组的方法,在原始方程组的单个或多个系数改变时,都需要从头开始重新执行整个过程.因此,如果对原始系统做出修改,之前这些计算所付出的大量计算努力就白费了.

然而,有一些替代方法(例如Sherman-Morrison方法和Woodbury方法)允许在对初始系统进行某些特殊类型的修改时,以较少的额外乘法和加法次数来计算新的解,所提方法可以直接应用于这种动态更新过程.事实上,我们在第7节中展示了当原始方程组中增加或删除某些方程和(或)变量时,如何以极少的额外计算开销来更新解.

总之,在本章中,我们表明第1章中描述的方法对于解决方程组最重要的问题非常有用,包括相容性分析,求解,或求某些选定变量的解.我们还提供了一种当系统中增加或删除方程和(或)未知数时更新线性方程组解的方法.该方法首先求解初始方程组,并保留与线性系统相关的矩阵行生成的线性空间的正交空间及其补空间的信息.此外,还会存储与每个方程相关的主元列信息,以及线性相关方程所依赖的基本独立方程的信息.利用这些信息,在上述四种情况的任何一种下,都可以对修改后的系统解进行更新.